\documentclass[11pt]{article}
\usepackage[T1]{fontenc}
\usepackage{lmodern,microtype,amsmath,amssymb,amsthm,booktabs}
\usepackage[letterpaper,margin=.92in,top=.84in,bottom=.86in]{geometry}
\usepackage{xcolor}
\definecolor{inkblue}{RGB}{28,63,91}
\usepackage[colorlinks,linkcolor=inkblue,citecolor=inkblue,urlcolor=inkblue,pdfauthor={OpenAI},pdftitle={Distinguishing Dirichlet characters and abelian number fields by low zeros}]{hyperref}
\usepackage{fancyhdr}
\pagestyle{fancy}\fancyhf{}\fancyhead[L]{\small\itshape Dirichlet characters and abelian fields}\fancyhead[R]{\small\itshape Finite zero windows}\fancyfoot[C]{\thepage}
\setlength{\headheight}{14pt}
\newtheorem{theorem}{Theorem}
\newtheorem{corollary}[theorem]{Corollary}
\newtheorem{lemma}[theorem]{Lemma}
\newtheorem{proposition}[theorem]{Proposition}
\newcommand{\C}{\mathbb C}\newcommand{\Q}{\mathbb Q}\newcommand{\R}{\mathbb R}
\newcommand{\abs}[1]{\left|#1\right|}\newcommand{\norm}[1]{\left\|#1\right\|}
\newcommand{\dd}{\,d}\DeclareMathOperator{\sinc}{sinc}
\DeclareMathOperator{\ord}{ord}
\newcommand{\calP}{\mathcal P}\newcommand{\wh}{\widehat}
\allowdisplaybreaks[2]\emergencystretch=1em
\begin{document}
\begin{center}
{\Large\bfseries\color{inkblue} Distinguishing Dirichlet characters and abelian\par
number fields by low zeros\par}
\vspace{.65em}
{\small Research note \enspace$\cdot$\enspace 26 November 2026}
\end{center}
\begin{abstract}
Using Dirichlet GRH, we prove that the strictly positive zero ordinates,
counted with multiplicity, up to height $C\log^2(2Q)/\log\log(3Q)$ determine a primitive Dirichlet
character of conductor at most $Q$, uniformly in its order. The
zero-multiplicity measures of any fixed number of distinct characters have full
linear rank in this range. For abelian number fields of absolute discriminant at
most $D$, a height of $C\log^2(2D)$ determines the field, uniformly in the degree.
Both statements extend to translated intervals and give lower bounds on the
total discrepancy of the multiplicities. The two estimates use different
features of the same explicit formula. A sieve bounds contamination from nearby
prime powers on average over prime centres. An $\ell^2$ estimate controls that
contamination for a sum of characters with bounded total log conductor.
\end{abstract}

\section{The finite data and what they determine}
For a primitive Dirichlet character $\chi$ set
\[
 m_\chi(t)=\ord_{s=1/2+it}L(s,\chi)\qquad(t\in\R).
\]
This is zero whenever the function is nonzero at the given point. Write $m_K$ for the corresponding function of
the Dedekind zeta function of a number field $K$. All zeros in this paper are
nontrivial. We use the Dirichlet zero-line theorem of \cite{lean}, including
the Riemann zeta function. An ordinate is recorded with its sign and
multiplicity. For a nonreal character the sign distinguishes the spectra of
$\chi$ and $\bar\chi$, which are reflections of one another. Put
\[
 I=I(\tau,H)=[\tau-H,\tau+H],\qquad
 J(F,\tau)=\log\log\bigl(e^e F(3+|\tau|)\bigr).
\]
Logarithms are natural. All constants below are absolute unless a dependence
is displayed. A sum over $t\in I$ ranges over the distinct ordinates in the
union of the spectra in question.

\begin{theorem}[Abelian fields]\label{thm:fields}
There are constants $C,c>0$ with the following property. Let $K$ and $L$ be
distinct finite abelian extensions of $\Q$ in a fixed algebraic closure, with
absolute discriminants at most $D\ge3$. Let $r$ be the cardinality of the
symmetric difference between their groups of primitive Dirichlet characters.
For every $\tau\in\R$ and
\[
 H\ge C\{\log^2(2D)+J(D,\tau)\}
\]
one has
\begin{equation}\label{eq:fieldmain}
 \sum_{t\in I}|m_K(t)-m_L(t)|
 \ \ge\ c H\sqrt r\,\frac{\log\log(3D)}{\log(2D)} .
\end{equation}
In particular the multiset for $|t|\le C'\log^2(2D)$ determines an abelian
field of absolute discriminant at most $D$ up to $\Q$-isomorphism. The
constant is independent of the field degrees.
\end{theorem}

\begin{theorem}[Linear independence in a finite window]\label{thm:family}
Let $\chi_1,\ldots,\chi_r$ be distinct primitive Dirichlet characters
of conductors at most $Q\ge3$. This includes the conductor-one character.
Take $X=C_0r^2\log^2(2Q)$ for a sufficiently large absolute $C_0$. For
every $v\in\C^r$ and every $\tau\in\R$,
\begin{align}
 H&\ge C_1\{rX/\log X+J(Q,\tau)\}
 \quad\Longrightarrow\nonumber\\[-2pt]
 \sum_{t\in I}\abs{\sum_{j=1}^r v_jm_{\chi_j}(t)}
 &\ \ge\ c H\,\frac{\log X}{\sqrt{rX}}\,\norm{v}_2 .
 \label{eq:familymain}
\end{align}
Consequently the matrix $(m_{\chi_j}(t))_{t\in I,\,1\le j\le r}$ has
column rank $r$. Some $r$ ordinates in $I$ already give an invertible minor.
\end{theorem}

Taking $r=2$ and $v=(1,-1)$ gives the character cutoff in the abstract.
Taking differences of multiplicities of factors also identifies any product
of primitive Dirichlet $L$-functions with at most $k$ distinct factors,
each of conductor at most $Q$, from the pooled zero multiset up to
$O_k(\log^2(2Q)/\log\log(3Q))$. Indeed the union of the factor lists of
two candidate products has size at most $2k$. The exponents of the factors
can be arbitrary positive integers. These are uniqueness statements from
exact finite data. A discrepancy in multiplicity can occur at an ordinate
shared by both functions.

\paragraph{Relation to the two repositories.}
The readable edition of \cite{lean} states its public interface in
\texttt{grh/src/Statements.lean}. In \texttt{Main.lean},
\texttt{LFunction\_eq\_zero\_re} places every zero in the open critical
strip on the critical line. The source derives a one-sided estimate for a
smoothed, quadratically twisted prime remainder by comparing lattice
covolumes. The analytic reduction excludes zeros to the right of the line.
Here the zero-line theorem is used both for the measured spectra and for
the primitive characters inducing pairwise quotients. Those quotient
characters supply the many arithmetic measurements needed below.

The 621-page ProveIt article \cite{article} merges 103 reports. Its section
on quadratic characters constructs compact tests that isolate one prime
power while suppressing the tail over the zeros. Corollary
\texttt{rec:cor:quadsep} gives the quadratic cutoff
$O(\log^2(2Q)\log\log(3Q))$. The open questions following that section ask
whether the last factor can be removed and how the cutoff changes for complex
characters. Theorem~\ref{thm:family} improves that upper bound by a factor
of order $(\log\log Q)^2$ and covers arbitrary orders. The
article also studies full rank and signed differences at large height for
fixed families. Its high-window theorem in \texttt{sections/zer.tex}
obtains triple-logarithmic widths, with the starting height depending on
the family. Our estimates give conductor dependence from height zero onward.

\paragraph{Earlier work and scope of the comparison.}
Banks \cite{banks} gives an unconditional comparison on positive-height
windows for distinct primitive characters of a common modulus.
Steuding \cite{steuding} studies global determination and unique
factorizations from zero multisets. Eid's deposited abstract \cite{eid}
describes polylogarithmic-height determination and stability for bounded
character orders by isolating a distinguishing prime power. We could verify
the abstract only, so we make no exponent comparison with that preprint.
The refinements proved here are the conductor-uniform scale in
Theorem~\ref{thm:family}, its quantitative rank bound, and the
degree-uniform discriminant estimate in Theorem~\ref{thm:fields}.
Historical priority for these formulations remains to be established.

\section{The explicit formula at positive logarithmic arguments}
For $g$ supported on a compact subinterval of $(0,\infty)$, write
$\wh g(t)=\int_\R g(u)e^{itu}\dd u$. Let $a\in\{0,1\}$ be the parity of
$\chi$ and $\delta_\chi$ be one for conductor one and zero otherwise.
For a sufficiently smooth, possibly complex $g$, the explicit formula reads
\begin{align}
 \sum_\gamma \wh g(\gamma)
 &=-\sum_{n\ge2}\frac{\Lambda(n)\chi(n)}{\sqrt n}g(\log n)
   -\int_0^\infty g(u)\frac{e^{-(a+1/2)u}}{1-e^{-2u}}\dd u \nonumber\\
 &\hspace{1.1em}+\delta_\chi\int_0^\infty
           g(u)(e^{u/2}+e^{-u/2})\dd u .\label{eq:weil}
\end{align}
Here $\Lambda(p^k)=\log p$ for $k\ge1$ and $\Lambda(n)=0$ otherwise. The sum on the left repeats each ordinate according to its multiplicity.
This is the positive-support specialization of the Dirichlet explicit
formula, as in \cite[Lemma 2.1]{liuye}. For the signs, put
$\Phi(s)=\int g(u)e^{(s-1/2)u}\dd u$. Integrating
$-(L'/L)(s,\chi)\Phi(s)$ on a line to the right of one gives the displayed
prime sum with positive sign by inversion. Shifting left, a zero
contributes $-\Phi(\rho)$ and a pole contributes $+\Phi(1)$.
Summing the values at $-a-2j$ gives the integral with denominator
$1-e^{-2u}$. For zeta the list starts at $-2$, which restores $\Phi(0)$.
Equivalently the general explicit formula has a reflected prime term
$g(-\log n)$ and a conductor term proportional to $g(0)$, both zero here.
Compactly supported piecewise smooth tests used below follow by smoothing,
since their Fourier transforms decay integrably against the zero count.

We use the uniform count, with multiplicity,
\begin{equation}\label{eq:count}
 \#\{\gamma:t\le\gamma\le t+1\}\ll\log(2q(|t|+3)).
\end{equation}
For conductor $q>1$ this follows, for example, by differencing the
symmetric zero count of Bennett, Martin, O'Bryant and Rechnitzer
\cite{count}, allowing a slightly larger interval at the endpoints.
The zeta case is the classical Riemann--von Mangoldt formula.
In particular $\sum_\gamma (1+\gamma^2)^{-1}\ll \log(2q)$.

\begin{lemma}[Prime Gram estimate]\label{lem:gram}
Fix a smooth compactly supported real weight $w$ on $(0,\infty)$.
For primitive characters of conductors $q_j,q_k$, define
\[
 G_{jk}(w)=\sum_p(\log p)w(p/X)\chi_j(p)\overline{\chi_k(p)} .
\]
For distinct primitive labels, and also on the diagonal, one has
\[
 G_{jk}(w)=\mathbf1_{j=k}X\int_0^\infty w(y)\dd y
          +O_w\!\left(\sqrt X\log(2q_jq_k)\right)\qquad (X\ge X_w).
\]
\end{lemma}
\begin{proof}
Let $\eta$ be the primitive character inducing
$\chi_j\bar\chi_k$ modulo $\mathrm{lcm}(q_j,q_k)$. Its conductor is at
most $q_jq_k$, and it is principal exactly for $j=k$, by uniqueness of
primitive induction. Apply \eqref{eq:weil} to $\eta$ with
$g(u)=e^{u/2}w(e^u/X)$. For large $X$ this has positive support, and
\[
 \wh g(t)=X^{1/2+it}\int_0^\infty w(y)y^{-1/2+it}\dd y
       \ll_w \sqrt X(1+t^2)^{-1}.
\]
Equation \eqref{eq:count} bounds the zero sum by $O_w(\sqrt X\log(2q_jq_k))$.
The integral from trivial zeros is $O_w(1)$, and the pole terms for
$\eta=1$ are $X\int w+O_w(1)$. Remove higher powers at cost $O_w(\sqrt X)$:
Chebyshev's bound $\sum_{p\le y}\log p\ll y$ bounds the squares by
$O_w(\sqrt X)$ and the other powers by $O_w(X^{1/3}\log X)$.
Passing from $\eta(p)$ to the product of the given characters at primes
costs at most $\|w\|_\infty\sum_{p\mid q_jq_k}\log p$.
\end{proof}

Fix once and for all $0\le w\le1$, smooth, supported in $(1,2)$, with
$\int w>0$. Write $A(n)=\sum_j v_j\chi_j(n)$,
$P_X=X/\log X$, and $\calP_X=\{p\text{ prime}:X\le p\le2X\}$.
For $X\ge C_0r^2\log^2(2Q)$ the lemma gives
\begin{equation}\label{eq:rms}
 \sum_{p\in\calP_X}|A(p)|^2\gg P_X\|v\|_2^2,\qquad
 \frac1{P_X}\sum_{p\in\calP_X}|A(p)|\gg\frac{\|v\|_2}{\sqrt r}.
\end{equation}
The quadratic-form error is
$O(\sqrt X\log(2Q)\|v\|_1^2)$, and $|A(p)|\le\sqrt r\|v\|_2$.
We also have $|\calP_X|\asymp P_X$ by the prime number theorem.

\section{A prime and its neighbours}
A compact window around $\log p$ can contain several prime powers. We compare
their total contribution to the unit weight at the prime $p$.
Define $\sinc u=\sin(u)/u$, with $\sinc 0=1$. For $B>0$ put
\begin{equation}\label{eq:leak}
 W_p(B)=\sum_{\substack{n\ne p\\ |\log(n/p)|\le1/2}}
  \frac{\Lambda(n)}{\log p}\sqrt{\frac pn}\,
  \min\{1,(B\log(n/p))^{-2}\}.
\end{equation}
Each $W_p(B)$ bounds the off-diagonal arithmetic error in a single character value
when the test has modulus at most $\sinc^2(B\log(n/p))$ on this support.

\begin{lemma}[Mean contribution from neighbouring prime powers]\label{lem:sieve}
Uniformly for all sufficiently large $X$ and $B>0$,
\[
 \frac1{P_X}\sum_{p\in\calP_X} W_p(B)\ll\frac{X}{B\log X}.
\]
\end{lemma}
\begin{proof}
On the support in \eqref{eq:leak}, $X e^{-1/2}\le n\le 2Xe^{1/2}$.
The prefactor is bounded absolutely, and the minimum is at most
$b(h)=\min(1,(C X/(Bh))^2)$, where $h=|n-p|\ge1$.
For two primes at nonzero distance $h$, the upper-bound sieve gives,
uniformly in $h$ in this range,
\begin{equation}\label{eq:pairs}
 \#\{p\in\calP_X:p\pm h\text{ prime in }[Xe^{-1/2},2Xe^{1/2}]\}
 \ll \frac{X}{\log^2X}\,s(h),\qquad
 s(h)=\prod_{\substack{\ell\mid h\\\ell>2\ {\rm prime}}}
       \frac{\ell-1}{\ell-2}.
\end{equation}
Here either sign may be fixed. We give details on uniformity. Odd $h$ gives
no such pair for large $X$. For even $h$, sieve out the classes $0$ and
$\mp h$ modulo primes $\ell\le y$, taking one class at $\ell\mid h$ and
two otherwise. Let $R=\sqrt X$, $y=R^{1/8}$, and let $\nu_\ell$ count the
excluded classes. The arithmetic large sieve
\cite[Corollary 1]{mv} bounds the surviving count by $O(X+R^2)/J_0$,
where
\[
 J_0=\sum_{\substack{d\le R,\ d\ {\rm squarefree}\\
                         \ell\mid d\Rightarrow\ell\le y}}
                 \prod_{\ell\mid d}\frac{\nu_\ell}{\ell-\nu_\ell}.
\]
All the sieving primes lie below both entries of every counted pair.
In the Euler product without $d\le R$, the prime $\ell$ is selected
independently with probability $\nu_\ell/\ell$. The expected $\log d$ is
at most $2\sum_{\ell\le y}\log\ell/\ell\le 2\log y+O(1)$.
Markov's inequality retains at least half the product for $d\le R$ for
large $X$. By Mertens' product formula, the reciprocal of this product
is $\ll s(h)/(\log y)^2$. For odd $\ell\nmid h$ the ratio of
$1-2/\ell$ to $(1-1/\ell)^2$ is $1-1/(\ell-1)^2$, and changing to
one excluded class multiplies $1-2/\ell$ by $(\ell-1)/(\ell-2)$.
This proves \eqref{eq:pairs}.

By expanding $s(h)$ over squarefree divisors,
$\sum_{h\le z}s(h)\le z\prod_{\ell>2}(1+1/(\ell(\ell-2)))\ll z$.
Summing by dyadic intervals gives
$\sum_{h\ge1}b(h)s(h)\ll X/B$, including $B>X$, when the stronger
bound $O((X/B)^2)$ applies up to an absolute factor. Thus prime
neighbours contribute $O(X^2/(B\log^2X))$ to the sum over all centres.
There are $O(\sqrt X)$ proper prime powers in the whole range.
For each, summing over all integer centres other than itself costs
$O(\sum_{h\ge1}b(h))=O(X/B)$. After division by $P_X$ their contribution
is $O(\sqrt X\log X/B)$, which is absorbed in the asserted bound.
\end{proof}

\section{Tests on translated zero windows}
Take $B=H/4$ and $m=\lfloor H/200\rfloor$, with the absolute lower
bound on $H$ chosen so that $m\ge4$. Set
\[
 w_m(u)=(1-4u^2)_+^m,\qquad
 g_p(u)= e^{-i\tau(u-\log p)}w_m(u-\log p)\sinc^2(B(u-\log p)),
 \qquad h_p=\wh g_p .
\]
Here $x_+=\max(x,0)$. The phase is chosen so that $g_p(\log p)=1$ at
every centre height $\tau$. Uniformly in real $t$,
\begin{equation}\label{eq:lip}
 |h_p(t)|\le \pi/B,\qquad |h'_p(t)|\le (\log p+1/2)\pi/B .
\end{equation}
Both follow by integration since $\int_\R\sinc^2(Bu)\dd u=\pi/B$.
A compact polynomial cutoff also controls the contribution from every
zero outside the measured window:
\begin{equation}\label{eq:tail}
 \sum_{|\gamma-\tau|>H}|h_p(\gamma)|
 \ll H e^{-c_0H}\log(2q(|\tau|+H+3)).
\end{equation}
To see this, integrate $m$ times by parts in $\wh w_m$. Boundary terms
vanish, and the $L^1$ norm of the $m$th derivative on the support is at most
$10^m m!$ by Cauchy's derivative estimate on complex circles of radius one
for the polynomial $(1-4z^2)^m$. Thus
$|\wh w_m(t)|\le(10m/|t|)^m$ for $t\ne0$.
The factor $\sinc^2(Bu)$ is an average of $e^{ivu}$ for $v$ in $[-2B,2B]$
(the sum of two independent uniforms on $[-B,B]$).
For $|t-\tau|\ge H$, this gives
$|h_p(t)|\le(20m/|t-\tau|)^m$.
Summing on dyadic annuli with \eqref{eq:count} proves \eqref{eq:tail};
the $j$th annulus contributes at most a constant times
$2^jH(10^{-1}2^{-j})^m\log(2q(|\tau|+2^{j+1}H+3))$.

For any complex coefficients $v_j$ define the real-ordinate multiplicity
combination $M(t)=\sum_jv_jm_{\chi_j}(t)$ and the normalized observation
\[
 Z_v(p)=-\frac{\sqrt p}{\log p}\sum_{t\in I}h_p(t)M(t).
\]
Formula \eqref{eq:weil} gives an actual matrix identity
\begin{equation}\label{eq:observation}
 Z_v(p)=A(p)+\sum_{n\ne p}\frac{\sqrt p}{\log p}
          \frac{\Lambda(n)}{\sqrt n}\,g_p(\log n)A(n)+E_v(p),
\end{equation}
where, if $q_j\le Q$, the following bound holds for each $p\in\calP_X$:
\begin{equation}\label{eq:errors}
 |E_v(p)|\ll
 \frac{\|v\|_1}{B\log X}
 +|v_{\rm triv}|\frac{X}{B\log X}
 +\|v\|_1\frac{\sqrt X}{\log X}
       H e^{-c_0H}\log(2Q(|\tau|+H+3)).
\end{equation}
The coefficient $v_{\rm triv}$ is zero if zeta is absent. On the support
of $g_p$, the trivial-zero kernel in \eqref{eq:weil} is $O(X^{-1/2})$ and
the zeta pole kernel is $O(\sqrt X)$. Their contributions follow from
$\int |g_p|\le\pi/B$. The remaining term is \eqref{eq:tail}.
In particular no knowledge of the conductor or parity of an unknown
character is required to form the observation.

\begin{proof}[Proof of Theorem~\ref{thm:family}]
The main term has mean absolute value
$P_X^{-1}\sum_p|A(p)|\gg\|v\|_2/\sqrt r$ by \eqref{eq:rms}.
Taking absolute values of the off-diagonal term in \eqref{eq:observation}
bounds its mean by
\[
 \|v\|_1\,P_X^{-1}\sum_{p\in\calP_X}W_p(B)
 \ll\sqrt r\|v\|_2\,\frac{X}{B\log X}.
\]
For $B\ge C r X/\log X$ this can be made less than a quarter of the main
lower bound by choosing $C$ large. The first two terms of \eqref{eq:errors}
are also smaller than a quarter with a sufficiently large constant.
For the last one, $H\gg rX/\log X$ absorbs every fixed power of
$r,X,H$ into $e^{c_0 H/4}$ (increasing the constant if needed).
Also $\log(2Q(|\tau|+H+3))\ll \exp(J(Q,\tau))+\log(H+3)$.
Taking $H\ge C J(Q,\tau)$ absorbs this factor, with room to spare.
It follows that some prime $p$ satisfies
$|Z_v(p)|\gg\|v\|_2/\sqrt r$.
On the other hand, \eqref{eq:lip} yields
\[
 |Z_v(p)|\ll \frac{\sqrt X}{H\log X}\sum_{t\in I}|M(t)|.
\]
Combining the bounds proves \eqref{eq:familymain}. Full rank and the
existence of the minor follow by finite-dimensional linear algebra.
\end{proof}

\begin{corollary}[Stability under partial matching]\label{cor:match}
Use the hypotheses of Theorem~\ref{thm:family} for two distinct characters,
with $r=2$. Match any submultiset of the first list in $I$ bijectively to
any submultiset of the second list in $I$. Let $U$ be the total number
of unmatched entries, and $S$ the sum of absolute displacements over all
matched pairs. Then
\[
 U+(\log(2X)+1)S\gg H\frac{\log X}{\sqrt X}.
\]
In particular every complete matching has total displacement
$\gg H/\log(2Q)$, whenever the cardinalities agree.
\end{corollary}
\begin{proof}
For $v=(1,-1)$ the proof gives a prime with $|Z_v(p)|\gg1$.
The contribution of an unmatched entry is bounded using the first
inequality of \eqref{eq:lip}, and that of a matched pair using the second.
\end{proof}

\section{Abelian fields with unrestricted degree}
For a finite abelian extension $K/\Q$, let $\mathcal X_K$ be its group of
Dirichlet characters, each represented by its unique primitive character.
The product in this group is taken with primitive induction. The standard
factorization and conductor--discriminant formula are
\begin{equation}\label{eq:class}
 \zeta_K(s)=\prod_{\chi\in\mathcal X_K} L(s,\chi),
 \qquad |\operatorname{disc}K|=\prod_{\chi\in\mathcal X_K} q_\chi .
\end{equation}
The group determines $K$: inside a common cyclotomic field it has as
annihilator exactly the subgroup fixing $K$ \cite[Chs.~3--4]{wash}.
Thus if $K\ne L$, cancellation of the common factors leaves $r\ge1$
distinct characters with coefficients $v_j\in\{1,-1\}$. The trivial
character cancels. Then
\begin{equation}\label{eq:budget}
 M(t)=m_K(t)-m_L(t),\qquad
 S_0=\sum_{j=1}^r\log(2q_j)\ll\log(2D),\qquad r\ll\log(2D).
\end{equation}
Indeed each surviving $q_j\ge3$ and $\sum_j\log q_j\le2\log D$.
Estimating every off-diagonal Gram entry by the largest conductor would
lose the benefit of this budget.

\begin{proof}[Proof of Theorem~\ref{thm:fields}]
Take $X=C_2\log^2(2D)$ with a sufficiently large absolute $C_2$. In
Lemma~\ref{lem:gram} the error after summing over all $j,k$ has size
\[
 \ll_w \sqrt X\sum_{j,k}\log(2q_jq_k)\ll_w \sqrt X\,rS_0 .
\]
For our fixed minorant on $(1,2)$ this is less than half the diagonal
main term $rX\int w$. For a second fixed smooth weight majorizing the
indicator of $[e^{-1/2},2e^{1/2}]$ the same argument gives an upper
bound. Write $\mathcal N_X=[e^{-1/2}X,2e^{1/2}X]\cap\mathbb N$.
As in \eqref{eq:rms}, these bounds give
\begin{equation}\label{eq:fieldgram}
 \sum_{p\in\calP_X}|A(p)|^2\gg rP_X,\qquad
 \sum_{p\in\mathcal N_X\ {\rm prime}} |A(p)|^2\ll rP_X .
\end{equation}
Define $b_n=(\Lambda(n)/\log X)\sqrt{X/n}\,A(n)$ for $n\in\mathcal N_X$.
Its prime part satisfies $\sum_{p}|b_p|^2\ll rP_X$. For higher powers,
$|A(n)|\le r$, $\Lambda(n)\ll\log X$, and their total von Mangoldt
weight is $O(\sqrt X)$ by the Chebyshev argument in Lemma~\ref{lem:gram}.
Hence
\begin{equation}\label{eq:columnnorm}
 \sum_{n\in\mathcal N_X}|b_n|^2
 \ll rP_X+\frac{r^2\sqrt X}{\log X}\ll rP_X ,
\end{equation}
using $r\ll\sqrt X$. This step controls prime powers even when several
characters agree on them.

The off-diagonal matrix with rows $p\in\calP_X$ and columns
$n\in\mathcal N_X$ has entries
\[
 T_{pn}=\mathbf1_{p\ne n}\sqrt{p/X}\,(\log X/\log p)g_p(\log n).
\]
For $n\ne p$, $|\log(n/p)|\gg |n-p|/X$, so
$|T_{pn}|\ll(X/B)^2|n-p|^{-2}$. Both absolute row sums and absolute
column sums are $O((X/B)^2)$, since the indices are integers without
repetition. Cauchy--Schwarz in each row followed by the column bound gives
$\|Tb\|_2\ll(X/B)^2\|b\|_2$ (Schur's test). Thus $B\ge K_0 X$ with
a large absolute $K_0$ makes the norm of $Tb$ at most a quarter of the
norm of the central vector $(A(p))_{p\in\calP_X}$, by
\eqref{eq:fieldgram}--\eqref{eq:columnnorm}.

In \eqref{eq:errors} one has $v_{\rm triv}=0$, $\|v\|_1=r$, each $q_j\le D$, and $r\ll\sqrt X$.
Once $H\ge C(X+J(D,\tau))$ the error at each prime is at most an
arbitrarily small absolute constant times $\sqrt r$, by the exponential
estimate used in the proof of Theorem~\ref{thm:family}. Since
$|\calP_X|\ll P_X$, its norm is at most another quarter of the central
norm. Identity \eqref{eq:observation} now yields
\[
 \sum_{p\in\calP_X}|Z_v(p)|^2\gg rP_X .
\]
For some prime $p$ this gives $|Z_v(p)|\gg\sqrt r$. By \eqref{eq:lip},
$\sum_{t\in I}|M(t)|\gg H\sqrt r\,\log X/\sqrt X$.
Substitution of $X=C_2\log^2(2D)$ proves \eqref{eq:fieldmain}.
\end{proof}

\paragraph{Positive ordinates alone.}
Both determination results hold using only strictly positive ordinates
below the indicated height. Choose $H$ of the asserted order for a centred
window and use the translated window $[H,3H]$, centred at $2H$, increasing
the absolute constant as needed. Its $J$ term obeys the same bound.
This observation applies even to a conjugate pair of complex characters.
Folding both halves onto their absolute values would give identical folded
lists for a conjugate pair.

\section{A pair of biquadratic fields}
Consider
\[
 K=\Q(\sqrt5,\sqrt{221}),\qquad L=\Q(\sqrt{13},\sqrt{85}).
\]
They are totally real of degree four and both have discriminant
$1105^2=1\,221\,025$. Their quadratic subfields have discriminants
$5,221,1105$ and $13,85,1105$. They share the factors $\zeta$ and
$L(s,\chi_{1105})$. Here $\chi_d$ denotes the real primitive character
of fundamental discriminant $d$. At $3$ the four remaining characters give
\[
 \chi_5(3)+\chi_{221}(3)-\chi_{13}(3)-\chi_{85}(3)=-4 .
\]
In particular the two zeta functions differ at their Euler factors even
though their gamma and discriminant factors agree.

\begin{proposition}[A low zero for the example]\label{prop:example}
There exists $1.20<t<1.23$ such that $\zeta_K(1/2+it)=0$.
Throughout $1.20\le t\le1.23$ one has $\zeta_L(1/2+it)\ne0$.
\end{proposition}
\begin{proof}
For these positive fundamental discriminants the quadratic root number
is $+1$. The completed $L$-function on the line is the real function
\[
 F_d(t)=4\int_0^\infty e^{u/2}\cos(tu)
       \sum_{n\ge1}\chi_d(n)\exp(-\pi n^2 e^{2u}/d)\dd u.
\]
This follows by splitting the Mellin integral of the quadratic theta
series at $1$, applying its theta transformation to the lower part,
and substituting $x=e^{2u}$. The gamma and exponential factors relating
$F_d(t)$ to $L(1/2+it,\chi_d)$ are nonzero. The following are outward
enclosures from \texttt{theta\_example.py}. In the last three rows the
enclosure is for the entire interval.
\begin{center}
\begin{tabular}{rll}
\toprule
$d$ & $t$ & enclosure for $F_d(t)$ \\
\midrule
221 & $1.20$ & $[0.040,\ 0.122]$ \\
221 & $1.23$ & $[-0.150,\ -0.067]$ \\
13  & $[1.20,1.23]$ & $[1.640,\ 1.693]$ \\
85  & $[1.20,1.23]$ & $[2.393,\ 2.656]$ \\
1105 & $[1.20,1.23]$ & $[3.727,\ 6.562]$ \\
\bottomrule
\end{tabular}
\end{center}
The checker integrates interval ranges on $2560$ subintervals of $[0,5]$,
each of length $1/512$, using endpoints on the $2^{-80}$ grid and outward
rounding. It encloses $\pi$ by Machin's identity, and exponentials and
cosines by Taylor estimates. At most $150$ terms are used in each theta
series. Terms among them dropped because the exponent exceeds $60$
contribute at most $150e^{-60}$ inside the sum. The remaining $n>150$
terms contribute at most $1105e^{-3\cdot150^2/1105}/900$ by Gaussian
integration. Multiply both bounds by $20e^{5/2}$ to account for the
integral on $[0,5]$. For $u=5+z$, the full theta sum has absolute value
at most $2e^{-40(1+2z)}$, so the omitted integral costs at most
$8e^{5/2-40}/(80-1/2)$. The program includes $10^{-9}$ for the total
of these errors.

The intermediate value theorem gives a zero of $F_{221}$ between the
endpoints, while $F_{13},F_{85},F_{1105}$ stay positive. It remains to
check the zeta factor of $L$. For $s=1/2+it$ Euler--Maclaurin gives
\[
 \zeta(s)=\frac1{s-1}+\frac12+\frac{s}{12}+R(s),\qquad
 |R(s)|\le\frac{|s(s+1)|}{12(1+\Re s)} .
\]
The inequality follows by bounding the second periodic Bernoulli polynomial
by $1/6$ in the integral remainder from $1$ to infinity. Over the stated
interval the imaginary part of the three displayed terms is
$-t/(1/4+t^2)+t/12<-0.57$, and $|R(s)|<0.16$. All four factors of
$\zeta_L$ are therefore nonzero there. This verification is independent
of GRH.
\end{proof}

\paragraph{A second certificate counts multiplicity differences.}
With the Dirichlet zero-line input, the same fields satisfy
$\sum_{0<t\le80}|m_K(t)-m_L(t)|\ge7$. Here the count allows several
units of discrepancy at one ordinate. The rational calculation in
\texttt{field\_example.py} is as follows. Put $a=2/7$,
$\phi(x)=(x(1-x^2)_+^7)'$, and $g(u)=\phi((u-\log3)/a)$. Its support
contains only $\log3$ among logarithms of prime powers. Its integral is
zero, and $g(\log3)=1$. In the four-character contrast of
\eqref{eq:weil} the trivial-zero integrals cancel by parity. The full
spectral sum is $4\log(3)/\sqrt3>263/105$. Six integrations by parts give
$|\wh g(t)|\le C_6 |t|^{-6}$ with
\[
 C_6=a^{-5}\left(2\int_{-1}^1|\phi^{(6)}(x)|^2\dd x\right)^{1/2}
 \le 2629707255/16 .
\]
Let $N(t)$ count both signs over the four characters of conductors
$5,13,85,221$ and set $A(t)=(4t/\pi)\log(\sqrt{1105}\,t/(2\pi e))$.
By \cite[Corollary 1.2]{count}, $|N(t)-A(t)|\le E(t)$ for $t\ge1$,
where $E(t)=(6/5)\log(34t)+18$. Partial summation bounds the absolute
tail beyond $H$ by
\[
 C_6\left(\int_H^\infty A'(t)t^{-6}\dd t+E(H)H^{-6}
                   +6\int_H^\infty E(t)t^{-7}\dd t\right).
\]
At $H=80$ this is less than $0.117$. The central ordinate contributes
zero. The two terms at $\pm t$ cost at most
$2\|g\|_1=8a(14/15)^7/\sqrt{15}$ per unit of positive-height
discrepancy. Subtraction of the tail leaves more than six times this amount.

\paragraph{A complex character modulo five.}
The script \texttt{prime\_leakage.py} sums all prime powers
in the support using $w_4(u)\sinc^2(Bu)$, which bounds the absolute
weight for all $m\ge4$. At $X=1000$, $B=4X/\log X$, $p=1523$, and
$\chi$ modulo five with $\chi(2)=i$, the central term is $-i$.
The normalized arithmetic sum is approximately
$0.04284475-1.02239786i$. Its distance from $-i$ is approximately $0.04834601$,
below the total absolute off-diagonal contribution $0.10965114$ (rounded
up). For $\bar\chi$ the arithmetic sum is its complex conjugate. The
JSON file gives all nine parameter choices. The prime sums use floating-point arithmetic. Proposition~\ref{prop:example} uses a
separate interval calculation.

\section{Further research questions}
\begin{enumerate}
\item Can the degree-uniform field cutoff be reduced by another
$\log\log D$? The Schur estimate treats possible neighbours as integers.
Improvement would require controlling the restriction of this matrix to the
particular prime-power vectors furnished by abelian character groups.
\item How many \emph{distinct} ordinates, as opposed to multiplicity
discrepancies, must distinguish two fields? The lower bounds here permit
concentration at a common zero of different orders. A uniform guarantee
of a zero present on only one side would have different arithmetic content.
\item Can the $J(D,\tau)$ term be reduced uniformly in the discriminant
while retaining the measure bound? The short-window results for fixed
families in \cite{article} suggest testing more rapidly decaying
spectral kernels against the prime-power estimates of this paper.
\item What absolute constant gives a practical degree-uniform cutoff?
The proof is effective in principle. Numerical optimization would require
explicit constants for the Gram bounds, the prime-power operator norm,
and the spectral tail simultaneously.
\end{enumerate}

\paragraph{Verification and provenance.}
The public Lean declarations and the authors' account of the machine checks
were inspected at the revisions below. The readable edition pins Lean
4.35.0 and Mathlib \texttt{56a00585}, and reports the axioms
\texttt{propext}, \texttt{Classical.choice}, and \texttt{Quot.sound}.
We did not rerun the upstream build. The Lean development supplies Dirichlet GRH and the ProveIt article gives
the comparison theorem and open question. The proofs in this note are
analytic arguments. The accompanying computations check the specified
finite examples.

\begin{thebibliography}{99}\small
\bibitem{lean} L. Sernau and collaborators,
\emph{GRH for Dirichlet L-functions, in Lean 4}, readable edition,
commit \texttt{72fead102f663d5c7a407964a13dc0cf9d977207}.
\href{https://github.com/lsernau/grh-lean/tree/72fead102f663d5c7a407964a13dc0cf9d977207/grh}{Repository and verification notes}.
\bibitem{article} V. Reshetnikov, ProveIt repository,
\emph{Beyond the Riemann hypothesis for Dirichlet L-functions: consequences,
transfers and limits}, commit \texttt{10802a6b360d5cdce48698261bc305e57a934b1f}.
\href{https://github.com/VladimirReshetnikov/ProveIt/tree/10802a6b360d5cdce48698261bc305e57a934b1f/Analysis/GeneralizedRiemannHypothesis/Article}{Article, source, and provenance};
in particular \href{https://github.com/VladimirReshetnikov/ProveIt/blob/10802a6b360d5cdce48698261bc305e57a934b1f/Analysis/GeneralizedRiemannHypothesis/Article/sections/rec.tex}{the section on recovery}.
\bibitem{banks} W. D. Banks, \emph{Comparing zeros of distinct Dirichlet
$L$-functions}, revised 2024, \href{https://arxiv.org/abs/2302.07073}{arXiv:2302.07073}.
\bibitem{steuding} J. Steuding, \emph{Distinguishing L-functions by joint
universality}, Lithuanian Math. J. \textbf{61} (2021), 413--423.
\href{https://doi.org/10.1007/s10986-021-09531-y}{DOI}.
\bibitem{eid} S. Eid, \emph{Effective Determination of Primitive Dirichlet
Characters by Their Zeros up to Polylogarithmic Height}, Zenodo, July 2026.
\href{https://doi.org/10.5281/zenodo.21326652}{Deposited record}.
\bibitem{liuye} J. Liu and Y. Ye, \emph{Distribution of zeros of Dirichlet
$L$-functions and the least prime in an arithmetic progression},
Acta Arith. \textbf{119} (2005), 13--38.
\href{https://doi.org/10.4064/aa119-1-2}{DOI}.
\bibitem{count} M. A. Bennett, G. Martin, K. O'Bryant, and A. Rechnitzer,
\emph{Counting zeros of Dirichlet $L$-functions}, Math. Comp. \textbf{90}
(2021), 1455--1482. \href{https://arxiv.org/abs/2005.02989}{arXiv:2005.02989}.
\bibitem{mv} H. L. Montgomery and R. C. Vaughan, \emph{The large sieve},
Mathematika \textbf{20} (1973), 119--134.
\href{https://doi.org/10.1112/S0025579300004708}{DOI}.
\bibitem{wash} L. C. Washington, \emph{Introduction to Cyclotomic Fields},
2nd ed., Graduate Texts in Mathematics 83, Springer, 1997, Chs.~3--4.
\href{https://doi.org/10.1007/978-1-4612-1934-7}{DOI}.
\end{thebibliography}
\end{document}
